\documentclass{article}
\usepackage[preprint]{neurips_2026}

\usepackage[utf8]{inputenc}
\usepackage[T1]{fontenc}
\usepackage{hyperref}
\usepackage{url}
\usepackage{booktabs}
\usepackage{amsfonts}
\usepackage{amsmath}
\usepackage{microtype}
\usepackage{xcolor}
\usepackage{graphicx}

\makeatletter
\renewcommand{\@noticestring}{ES 678 Course Project}
\makeatother

% YOUR-USERNAME
\newcommand{\repourl}{https://github.com/krdivyanshu343-pixel/ES678_Machine_Learning_Project_fall_2026}

\title{Next-Day PM2.5 Forecasting in Indian Cities: A Careful Comparison with Persistence}

\author{%
  Divyanshu Kumar \\
  Roll Number: 26510041 \\
  ES 678 \\
  Indian Institute of Technology Gandhinagar
}

\begin{document}
\maketitle

\begin{abstract}
We ask whether today's pollutant readings and calendar information predict tomorrow's city-level PM2.5 in Indian cities better than persistence (tomorrow $=$ today). Using daily CPCB data for 17 cities, we train on 2015--2018, tune on 2019 and test once on January--July 2020, with leak-free features and a chronological split. Ridge, random forest and gradient boosting are compared with persistence and a city-month mean. The best model (random forest) only matches persistence on the test set (MAE 10.98 vs.\ 11.03\,$\mu$g/m$^3$; skill $+0.004$, 95\% block-bootstrap CI $[-0.018,+0.030]$). Lag features help, peaks are under-predicted, and split-conformal intervals reach 90\% coverage overall (0.93) but only 0.45 on high-pollution days.
\end{abstract}

\section{Problem and Motivation}
PM2.5 is the pollutant most linked to health damage in Indian cities, and a reliable next-day forecast lets people and authorities act early (alerts, school and traffic decisions). \emph{Question:} can information available at the end of day $t$ (pollutant readings, recent PM2.5 history, calendar) predict PM2.5 on day $t{+}1$ in the same city more accurately than persistence, $\hat y_{c,t+1}=y_{c,t}$? One row is one city on one day; the target $y_{c,t+1}$ is a regression target in $\mu$g/m$^3$, and a prediction is made at the end of day $t$ using data up to $t$ only. The claim is limited to predictive association for these cities and period; it says nothing about the causes of pollution. The test period (January--July 2020) lacks the October--December winter peak and includes India's COVID-19 lockdown (from 25 March 2020), a distribution shift we analyse explicitly.

\section{Data and Preparation}
We use the public CPCB daily air-quality data for 26 Indian cities (2015-01-01 to 2020-07-01; 29,531 rows) in the Kaggle release ``Air Quality Data in India'' \citep{rao}. Variables are PM2.5, PM10, NO, NO$_2$, NO$_x$, NH$_3$, CO, SO$_2$, O$_3$, benzene, toluene, xylene and AQI; the data are public and contain no personal information. The AQI column is never used, because it is computed from the same pollutants. \emph{Quality checks:} there are no duplicate city-days; PM2.5 is missing in 15.6\% of rows (xylene 61\%, PM10 38\%, NH$_3$ 35\%); PM2.5 ranges from 0.04 to 950 $\mu$g/m$^3$ (median 48.6). Coverage is uneven: nine cities have no PM2.5 reading before 2019. We keep the 17 cities with at least 200 usable training rows (dropped: Aizawl, Bhopal, Chandigarh, Coimbatore, Ernakulam, Guwahati, Jorapokhar, Kochi, Shillong). A row is usable if PM2.5 is observed on days $t$ and $t{+}1$, giving 13,124 training, 6,040 validation and 3,013 test rows.

\emph{Preprocessing.} Each city is reindexed to a complete daily calendar and gaps are filled by a past-only forward fill (at most 3 days). Features are the day-$t$ pollutants, PM2.5 lags 1--7, trailing 3- and 7-day means, cyclic month and weekday, and the city. Every learned step (median imputation, scaling for ridge, one-hot encoding, dropping features missing in more than 50\% of training rows, which removes only xylene) is fit on the training years only. A unit test confirms that changing data after day $t$ leaves the day-$t$ features unchanged.

\section{Methods}
\emph{Baselines.} Persistence, and a city-month mean (mean training target of the same city and calendar month). \emph{Models.} Ridge regression minimises $\|y-Xw\|_2^2+\lambda\|w\|_2^2$ and is stable under correlated pollutants \citep{hoerl}. A random forest averages 300 decorrelated regression trees and captures non-linear effects without scaling \citep{breiman}. Gradient boosting adds shallow trees fit stagewise to the gradient of the loss \citep{friedman}. All are implemented in scikit-learn \citep{sklearn} and trained with squared error; we report MAE, so the mismatch (squared loss targets the conditional mean, which exceeds the median for skewed PM2.5) is a known limitation. Complexity was added only if validation error improved. Tuned on 2019: ridge $\lambda$, forest depth and leaf size, boosting trees, rate and depth.

\emph{Intervals} (Section~\ref{sec:depth}). Quantile boosting fits the 5\% and 95\% conditional quantiles \citep{koenker}. Split-conformal prediction \citep{lei} returns $\hat y\pm\hat q$, where $\hat q$ is the $\lceil (n{+}1)(1{-}\alpha)\rceil$-th smallest absolute residual on a calibration year ($\alpha=0.1$). Its guarantee needs exchangeability, which time series do not satisfy \citep{barber}, so we measure coverage on later periods instead of assuming it.

\section{Evaluation Design}
The split is chronological by target date: train 2015--2018, validation 2019 (hyperparameters, model selection and conformal calibration), test 1 January--1 July 2020, evaluated once after all choices were frozen. A rolling-origin check (train $\le$2016, validate 2017; train $\le$2017, validate 2018) tests stability. Metrics: MAE (typical error, robust to outliers), RMSE (penalises large misses), and MAE on \emph{high-pollution days} (target above the 90th percentile of training targets, 172.6\,$\mu$g/m$^3$). The skill score is $1-\mathrm{MAE}_{\text{model}}/\mathrm{MAE}_{\text{persistence}}$; positive means better than persistence. Because consecutive days are correlated, 95\% intervals come from a moving-block bootstrap \citep{kunsch} over test dates (7-day blocks, all cities of a date resampled together, 2000 draws). Test results are reported overall, before the lockdown (to 24 March) and during it.

\section{Results and Error Analysis}
\begin{table}[t]
\caption{2020 test results ($n=3013$; MAE and RMSE in $\mu$g/m$^3$; high $=$ 38 days above 172.6; pre/lock $=$ before/during lockdown). Skill is relative to persistence with block-bootstrap 95\% CI.}
\label{tab:main}
\centering\small\setlength{\tabcolsep}{3pt}
\input{tables/table_main}
\end{table}
Random forest had the lowest validation MAE (13.38 vs.\ 13.63 for persistence) and was therefore selected. On the test set (Table~\ref{tab:main}) it ties persistence: MAE 10.98 vs.\ 11.03, skill $+0.004$ with a CI that includes zero (paired MAE difference $[-0.35,+0.20]$\,$\mu$g/m$^3$). Its RMSE is lower (17.5 vs.\ 18.4), so it makes somewhat fewer large misses. Gradient boosting is indistinguishable from persistence and ridge is clearly worse. The city-month mean fails (MAE 34.8, bias $+31.6$), because 2020 levels lie far below earlier years. The pattern is stable: forest skill is $+0.046$ and $+0.007$ in the 2017 and 2018 rolling folds and $+0.018$ on 2019. It beats persistence in 10 of 17 test cities, with per-city skill between $-0.09$ and $+0.10$. Before and during the lockdown its skill is $+0.014$ $[-0.019,+0.046]$ and $-0.011$ $[-0.037,+0.027]$; absolute errors are far smaller during the lockdown (7.7 vs.\ 14.8) because levels fell, not because the model improved. On the 38 high-pollution days MAE is 46.2 vs.\ 50.7 (skill CI $[-0.10,+0.22]$), too few days for a firm claim.

\begin{figure}[t]
\centering
\includegraphics[width=\linewidth]{figures/fig_errors.pdf}
\caption{(a) Observed vs.\ predicted PM2.5 on the 2020 test set; red points are high-pollution days. (b) MAE by target month. Errors are largest in January and are shared with persistence.}
\label{fig:errors}
\end{figure}
Representative errors (Figure~\ref{fig:errors}) are abrupt changes that history cannot anticipate: Delhi on 1 January 2020 (observed 372, predicted 233) and sudden drops in Brajrajnagar during the lockdown. Persistence misses the same days.

\section{Graduate Depth Analysis}\label{sec:depth}
We ran three analyses on this same task; each model family is re-tuned on 2019 for every feature set.

\emph{(1) Lag ablation.} Table~\ref{tab:abl} shows that history matters: for the forest, test MAE falls from 12.1 (today only) to 11.4 (lags 1--3) to 11.0 (lags 1--7), and validation MAE shows the same order (14.03, 13.58, 13.38); ridge and boosting agree. Without history, every model is worse than persistence, because squared-loss models shrink toward the mean while persistence copies today's level. Even with the full set the model only reaches persistence level, so in this test the extra pollutants and calendar terms do not lift it clearly above persistence.

\begin{table}[t]
\caption{Lag ablation: test MAE ($\mu$g/m$^3$) and skill vs.\ persistence in parentheses.}
\label{tab:abl}
\centering\small\setlength{\tabcolsep}{4pt}
\input{tables/table_ablation}
\end{table}
\emph{(2) Failure audit.} On the 38 high-pollution test days the forest under-predicts 71\% of the time (mean bias $-28.6$\,$\mu$g/m$^3$, predictions 87\% of observed); on the 266 high validation days the figures are 74\% and 87\%. Persistence also under-predicts (66\%, bias $-18.3$), partly mechanically: selecting days by a high outcome selects days that followed lower values, so any predictor looks low. The forest nevertheless shrinks peaks more than persistence does. In winter episodes it is only marginally better (Nov--Dec 2019 validation MAE 20.1 vs.\ 20.4; Jan--Feb 2020 test 16.2 vs.\ 16.6).

\emph{(3) Uncertainty and calibration.} Table~\ref{tab:int} gives coverage of nominal 90\% intervals; calibration uses the year before each test window. Split-conformal coverage is 0.95, 0.92 and 0.93 overall (2020 CI $[0.91,0.95]$) but unstable across regimes (0.88 before and 0.98 during the lockdown, since calibration residuals from 2019 are larger than lockdown errors), and its fixed width ($\pm30.7$) cannot follow peaks: coverage on high days is only 0.45--0.72. Quantile boosting under-covers slightly overall (0.87--0.88, 2020 CI $[0.86,0.89]$) but widens on peaks (mean width 166 vs.\ 44) and covers 0.74--0.85 of high days. Marginal coverage therefore hides failures on the days that matter most for alerts.

\begin{table}[t]
\caption{Interval coverage (target 0.90) and mean width ($\mu$g/m$^3$); the last columns give coverage on high-pollution days of the whole window.}
\label{tab:int}
\centering\footnotesize\setlength{\tabcolsep}{3pt}
\input{tables/table_intervals}
\end{table}

\section{Limitations and Conclusion}
The main limitation is the test window: six months, no October--December peak, one lockdown regime, and only 38 high-pollution days, so conclusions about peaks have wide uncertainty. City coverage is uneven (206--1457 training rows per city), station sets change over time, and models use squared loss. A stronger claim needs a multi-year test that includes winter episodes, station-level data, meteorology (not in this file), and loss functions aimed at MAE or quantiles. Predictive association is not causation. \emph{Answer:} with these data, learned models do not beat persistence by a meaningful margin (best skill $+0.004$, CI containing zero); recent PM2.5 history carries the signal, peaks are under-predicted, and split-conformal intervals are well calibrated on average but not on high-pollution days.

\section{Reproducibility and AI-Use Disclosure}
Code, notebook, results and LaTeX source: \url{\repourl}. The data need a free Kaggle login and are not redistributed; the README gives download steps. Splits, seeds (42), library versions and tuned settings are in \texttt{results/run\_summary.json}; \texttt{python scripts/run\_all.py} regenerates every number, table and figure (about 10 minutes on one CPU core). \emph{AI use:} Claude (Anthropic) assisted with planning, writing and debugging the code and LaTeX, and drafting this text. Outputs were checked with automated leakage and split tests, recomputation of the headline metrics from saved predictions, and inspection of the largest errors; the author is responsible for all content. No private data, passwords or API keys are used.

\begin{thebibliography}{9}
\bibitem[Rao(2020)]{rao} R.~Rao. Air Quality Data in India (2015--2020). Kaggle dataset, data made public by the Central Pollution Control Board, India. \url{https://www.kaggle.com/datasets/rohanrao/air-quality-data-in-india}
\bibitem[Pedregosa et~al.(2011)]{sklearn} F.~Pedregosa et al. Scikit-learn: Machine learning in Python. \emph{JMLR} 12:2825--2830, 2011.
\bibitem[Hoerl and Kennard(1970)]{hoerl} A.~E. Hoerl and R.~W. Kennard. Ridge regression: Biased estimation for nonorthogonal problems. \emph{Technometrics} 12(1):55--67, 1970.
\bibitem[Breiman(2001)]{breiman} L.~Breiman. Random forests. \emph{Machine Learning} 45(1):5--32, 2001.
\bibitem[Friedman(2001)]{friedman} J.~H. Friedman. Greedy function approximation: A gradient boosting machine. \emph{Annals of Statistics} 29(5):1189--1232, 2001.
\bibitem[Koenker and Bassett(1978)]{koenker} R.~Koenker and G.~Bassett. Regression quantiles. \emph{Econometrica} 46(1):33--50, 1978.
\bibitem[Lei et~al.(2018)]{lei} J.~Lei et al. Distribution-free predictive inference for regression. \emph{JASA} 113(523):1094--1111, 2018.
\bibitem[Barber et~al.(2023)]{barber} R.~F. Barber, E.~J. Cand\`es, A.~Ramdas and R.~J. Tibshirani. Conformal prediction beyond exchangeability. \emph{Annals of Statistics} 51(2):816--845, 2023.
\bibitem[K\"unsch(1989)]{kunsch} H.~R. K\"unsch. The jackknife and the bootstrap for general stationary observations. \emph{Annals of Statistics} 17(3):1217--1241, 1989.
\end{thebibliography}

\end{document}
